% Source: https://artofproblemsolving.com/community/c3h1411369_a_limerick
% Stable identity: aops-topic-1411369-q01
Let the age be $a$. We need $\dfrac{2}{3}a^2$ to be a perfect cube (an integer cube). Write $a=2^p\,3^q\,t$ with $\gcd(t,6)=1$. Then
\[
\dfrac{2}{3}a^2=2^{2p+1}\,3^{2q-1}\,t^2 .
\]
This is an integer only if $q\ge1$, and it is a cube exactly when every prime exponent is a multiple of $3$:
\[
2p+1\equiv0\pmod3\iff p\equiv1\pmod3,\qquad 2q-1\equiv0\pmod3\iff q\equiv2\pmod3,
\]
and $t^2$ is a cube, so $t$ is a cube. Thus $p\in\{1,4,7,\dots\}$, $q\in\{2,5,\dots\}$ and $t\in\{1,125,343,\dots\}$. The smallest values of $a=2^p3^qt$ are
\[
2\cdot9=18,\qquad 2^4\cdot9=144,\qquad 2\cdot3^5=486,\ \dots
\]
so the two smallest ages are $\boxed{18\text{ and }144}$. Check: $\dfrac{2}{3}\cdot18^2=216=6^3$ and $\dfrac{2}{3}\cdot144^2=13824=24^3$. (Only $18$ is a realistic age for a young lady.)
